%----------------------------------------------------------------------------------------
% Introduction
%----------------------------------------------------------------------------------------
\setcounter{page}{1} % Sets counter of page to 1
\renewcommand\thesection{\Roman{section}}
\renewcommand\thesubsection{\Alph{subsection}}
\section{Introduction} % Add a section title
\subsection{Background on Stablecoins} % Add a subsection
Cryptocurrency's promise of decentralized finance is often overshadowed by its inherent price volatility, hindering its everyday usability. Stablecoins emerged as a solution, offering the benefits of digital currencies with the stability of traditional assets. They provide a reliable medium of exchange, a store of value, and a unit of account suitable for daily transactions and e-commerce. 
However, this model has several drawbacks, including centralization, counterparty risk, and lack of transparency.\newline
USB, a crypto-asset-backed stablecoin, represents an alternative to traditional fiat-backed stablecoins. 
It uses a disruptive way of pegging to the US dollar by shorting crypto assets, using an inverse perpetual swap. 
This mechanism offers several benefits, including decentralization and enhanced privacy for users. 
This whitepaper provides an overview of USB, its architecture, mechanism, and future outlook.
 

\subsection{Problem with Fiat-backed Stablecoins} % Add another subsection
Fiat-backed stablecoins, such as Tether (USDT) and USD Coin (USDC), rely on a reserve of fiat currency held in a bank account to maintain their peg to the US dollar.
While this design has been successful in providing a stable digital asset, it also has significant drawbacks.\newline
Firstly, fiat-backed stablecoins are centralized and rely on trusted third parties, such as banks, to hold and manage the reserve.
This centralization introduces counterparty risk, as users must trust that the third party is holding the reserve in full and will honor redemption requests.\newline
Secondly, fiat-backed stablecoins are subject to regulatory scrutiny and can face crackdowns or frozen funds, limiting their availability and stifling creativity in the ecosystem.
This centralization also limits the privacy of users, as the fiat reserve and associated transactions are subject to regulatory oversight.\newline
Lastly, the reliance on fiat currency as a backing asset limits the scalability of fiat-backed stablecoins.
As the demand for stablecoins increases, so does the need for the reserve currency to back them, leading to potential liquidity issues and limiting the growth of the ecosystem.\newline
These limitations and risks associated with fiat-backed stablecoins have created a need for an alternative stablecoin solution that is decentralized and scalable.
This is where USB, the crypto-asset-backed stablecoin developed by Stabolut, comes in.


\subsection{Solution: Crypto-asset-backed Stablecoins like USB}
In order to address the limitations of fiat-backed stablecoins, we propose the use of crypto-asset-backed stablecoins like USB.
USB is a stablecoin that is backed by a basket of crypto assets and is designed to mimic the value of the US dollar.\newline
USB uses a disruptive mechanism of shorting crypto assets through an inverse perpetual swap to maintain its peg to the US dollar.
This allows USB to be free from the constraints of traditional fiat-backed stablecoins.\newline
Furthermore, USB does not require bank accounts or other centralized entities, making it more decentralized and less prone to regulatory crackdowns or frozen funds.
This provides users with increased privacy and security, as well as reduced counterparty risk.\newline
Overall, crypto-asset-backed stablecoins like USB represent a powerful solution to the limitations of fiat-backed stablecoins, offering greater transparency, decentralization, and security.


\subsection{Regulatory Alignment}  
USB’s design adheres to global regulatory frameworks by:  
\begin{itemize}  
\item Avoiding direct yield on USB (rewards are in SBL, a separate utility token).  
\item Structuring treasury airdrops as non-recurring, non-guaranteed distributions (avoiding classification as securities).  
\item Partnering with legal advisors to ensure jurisdictional compliance.  
\end{itemize}  

\subsection{Advantages of USB Over Other Stablecoins}
USB offers several advantages over traditional fiat-backed stablecoins and other crypto-asset-backed stablecoins, including:
\begin{itemize}
\item Decentralization: USB is a decentralized stablecoin, meaning that it is not controlled by any centralized authority, such as a bank or government.
\item Transparency: USB is transparent, with real-time auditing of the assets that back the stablecoin. This ensures that users can trust that the stablecoin is fully backed by assets.
\item Regulatory Alignment: USB staking rewards SBL tokens, avoiding legal conflicts in restrictive jurisdictions.
\item SBL Utility: The token accrues value from treasury airdrops, creating a sustainable ecosystem for holders.
\item Security: USB's security is based on its decentralized nature and the security of the underlying blockchain networks.
\item No KYC: USB does not require any Know Your Customer (KYC) checks, making it easier for users to access the stablecoin.
\item Lower counterparty risk: USB's asset backing is secured by a decentralized network of smart contracts, which reduces counterparty risk.
\item Passive Income Generation: USB's hedging strategies generated 60\% (ETHUSD) and 16\% (BTCUSD) yields in 2024 backtesting, outperforming traditional stablecoins. It is important to note that past performance is not indicative of future results, and these figures should not be considered a guarantee of future returns.
\item Minimal price volatility: USB's peg to the US dollar is maintained by shorting crypto assets using inverse perpetual swaps, which helps to minimize price volatility.
\item Overall, USB is a more secure, transparent, and decentralized stablecoin compared to traditional fiat-backed stablecoins, making it an ideal choice for users looking for a stablecoin with minimal risk and maximum security.
\end{itemize}
